\section{生成的本质}

\subsection{生成即采样}

我们先看几类典型的数据对象：
\begin{itemize}
  \item 图像：若图像大小为 $H\times W$，并带有 RGB 三个通道，则一张图像可以看成 $\vz \in \R^{H\times W\times 3}$；
  \item 视频：如果一段视频有 $T$ 帧，那么它可以看成 $\vz \in \R^{T\times H\times W\times 3}$；
\end{itemize}

\emph{因此，我们可以把待生成对象统一表示成一个向量 $\vz \in \R^d$。}


\begin{definition}[对象的向量表示]
在连续型生成建模中，我们把待生成对象抽象为高维向量 $\vz \in \R^d$，并把生成任务理解为在该空间中的概率建模与采样问题。
\end{definition}

假设我们想“生成一张狗的图片”。直觉上，这并不是要求模型输出某一张唯一正确的图像，因为满足要求的图像其实有很多张。更合理的理解是：我们希望模型能够产生\emph{一类合理样本}，而不是某个唯一答案。

在概率论视角下，这意味着我们引入一个数据分布 $\data$。它描述在对象空间 $\R^d$ 中，哪些样本更“像真实数据”，哪些样本不太合理。于是，“生成”就不再被理解为构造一个固定答案，而是理解为从这个数据分布中采样。

\begin{definition}[生成建模问题]
设 $\data$ 是定义在 $\R^d$ 上的数据分布。生成建模的目标是构造一个模型，使其能够输出近似服从 $\data$ 的样本，即 $\hat{\vz} \sim p_{\theta}$，$ p_{\theta} \approx \data$。
\end{definition}

这里 $p_{\theta}$ 表示模型所诱导的分布，参数 $\theta$ 则对应模型需要学习的部分。于是，生成模型本质上就是一个\emph{采样器}：训练完成后，我们希望它能够不断输出看起来像真实数据的样本。

\subsection{数据集的意义}

上面的定义里有一个关键难点：真实的数据分布 $\data$ 往往是未知的。我们并不能直接写出它的解析表达式，通常只能拿到有限个从中抽取出来的样本。而这部分样本则作为真实数据分布的代理。

\begin{definition}[经验样本]
生成建模中的数据集由有限个样本组成：$\vz_1,\vz_2,\ldots,\vz_N \sim \data$，这些样本共同作为真实数据分布 $\data$ 的有限观测。
\end{definition}

因此，训练生成模型并不是“记住所有样本”，而是利用这些样本去学习一个能够\emph{泛化地}近似 $\data$ 的采样机制。


\begin{summarybox}
\begin{enumerate}
  \item 待生成对象可以统一表示成高维向量；
  \item 生成任务可以形式化为从数据分布中采样；
  \item 训练时所用的数据集样本作为分布代理，帮助模型学习真实分布；
\end{enumerate}

\end{summarybox}
